\section{The tame quartic gate as a finite word evaluation}

Put $G=\widehat{\mathrm{GT}}$ and
\[
 D=\ker\lambda\cap\bigcap_{q\ \mathrm{prime}}\ker\rho_q^I.
\]
Write $g=(1,f_g)$ for an element of $D$, and use the signed
matrices
\[
 U=\begin{pmatrix}1&1\\0&1\end{pmatrix},\qquad
 L=\begin{pmatrix}1&0\\-1&1\end{pmatrix},\qquad
 S=ULU=\begin{pmatrix}0&1\\-1&0\end{pmatrix}.
\]
For an odd prime $\ell$, a subscript $\ell$ on these matrices
will mean reduction in $\mathrm{SL}_2(\mathbb F_\ell)$.
Let
\[
 \chi_{-4}(\ell)=\left(\frac{-1}{\ell}\right),\qquad
 m_\ell=\frac{\ell-\chi_{-4}(\ell)}2
 =\begin{cases}(\ell-1)/2,&\ell\equiv1\pmod4,\\
               (\ell+1)/2,&\ell\equiv3\pmod4.
   \end{cases}
\]

\begin{proposition}\label{prop:tame-harmonic-word-gate}
For every $g\in D$ the corrected harmonic frame satisfies
\begin{equation}\label{eq:tame-harmonic-1}
W_g=f_g(U^2,S)\qquad\text{in }\mathrm{SL}_2(\widehat{\mathbb Z}).
\end{equation}
Consequently its residue belongs to the cyclic group
\[
 \mathcal T_\ell
 =\{aI+bS_\ell:a,b\in\mathbb F_\ell,\ a^2+b^2=1\},
 \qquad |\mathcal T_\ell|=\ell-\chi_{-4}(\ell).
\]
The simultaneous compatibility of the radical and CM actions
on every rational quartic root is equivalent to the following
family of finite word identities:
\begin{equation}\label{eq:tame-harmonic-2}
\boxed{\quad
 f_g(U_\ell^2,S_\ell)^{m_\ell}=I
 \qquad(g\in D,\ \ell\ge5\text{ prime}).\quad}
\end{equation}
For a fixed $g$, the same equivalence holds pointwise. This is
a reformulation of the remaining tame quartic gate; it does
not assert the identities in \eqref{eq:tame-harmonic-2}.
\end{proposition}
\begin{proof}
The harmonic frame is $W_g=q(\beta_g)^{-1}D_g$.
On $D$, the established signed normalization gives $D_g=I$.
Proposition~\ref{prop:normalized-GF1-from-dyadic} gives
\[
 \beta_g(u^2,v^2)^{-1}=f_g(u^2,uvu)
\]
in the signed profinite braid group. Applying its continuous
signed representation $q(u)=U$, $q(v)=L$ proves \eqref{eq:tame-harmonic-1}, and
evaluation of a profinite word commutes with reduction modulo
$\ell$.

The frame belongs to the integral norm-one torus of
$\mathbb Z[i]$, represented by $i\mapsto S$.
Reduction gives the displayed group $\mathcal T_\ell$.
When $\ell\equiv1\pmod4$, its two eigenspaces identify it
with $\mathbb F_\ell^\times$; when $\ell\equiv3\pmod4$,
it is the kernel of the norm
$\mathbb F_{\ell^2}^\times\to\mathbb F_\ell^\times$.
Thus it is cyclic of the stated even order in both cases.
Its square subgroup is exactly the kernel of the power
$z\mapsto z^{m_\ell}$. The kernel of reduction from the
integral local torus is pro-$\ell$, so squaring is an
automorphism on that kernel. Hence a local frame component
at odd $\ell$ is a square if and only if its residue passes
the displayed power test.

Lemma~\ref{lem:lr-quartic-cm-quotient} identifies the
maximal elementary-two subextension of
$\mathbb Q(i)^{\mathrm{ab}}/\mathbb Q^{\mathrm{ab}}$
with
\[
 \mathbb Q^{\mathrm{ab}}(q^{1/4}:q\text{ rational prime}).
\]
Accordingly the action on all these quartic roots is trivial
if and only if $W_g$ is a square in the normalized adelic norm-one
group $N_4(4)$. The radical action is already trivial for $g\in D$.
At two the proved prime-two compatibility places $W_{g,2}$
in the local square subgroup. At three the established
harmonic residue is $W_g\bmod3\in\{I,-I\}$, which is
the square subgroup of the residue torus of order four;
the pro-three reduction kernel contributes no obstruction.
Thus precisely the primes $\ell\ge5$ remain.

Finally, squaring in the product defining $N_4(4)$ has image
equal to the product of their square subgroups: local square
roots may be chosen independently, and their tuple is an
adelic element of the required product. This retains the
simultaneous condition, rather than only checking a separate
list of necessary conditions. Substitution of \eqref{eq:tame-harmonic-1} now proves
both the pointwise and the universal equivalences in \eqref{eq:tame-harmonic-2}.
\end{proof}
